% GENERATED FILE. Edit canonical lessons, then run npm run generate:books.
% canonical-source-hash: fnv1a64:5663e06272de6243
% canonical-lessons: RU-C55-tree, RU-C55-flower, RU-C55-leaf, RU-C55-grass, RU-C55-village

\chapter{Tree, Flower, Leaf, Grass, Village}
\label{ch:ru-tree-flower-leaf-grass-village}
\hlchaptermodality{55}

\begin{chapteropening}
\textbf{By the end of this chapter:} \emph{I can name a tree, a flower, a leaf, grass and a village.}

Complete the last lesson of chapter 55: I can name a tree, a flower, a leaf, grass and a village.
\end{chapteropening}

\section[\texorpdfstring{\ru{дерево} (dérevo)}{дерево (dérevo)}]{\ru{дерево} (dérevo) — a tree}

\label{lesson:RU-C55-tree}



\begin{quote}
Before the new one: say the Russian for a ticket, then the Russian for a map.
\end{quote}

\subsection*{You'll want to know: \ru{дерево}}
\textbf{\ru{дерево}} (\emph{dérevo}) — \textquotedblleft{}a tree\textquotedblright{}.

\begin{cousinweb}
\textbf{\ru{дерево}} $\leftarrow$ Proto-Slavic \emph{*dervo} — inherited. Vasmer lists Lithuanian \emph{dervà}, \textquotedblleft{}pine\textquotedblright{}, Greek \emph{dóry}, \textquotedblleft{}wood, spear\textquotedblright{}, \emph{drŷs}, \textquotedblleft{}oak\textquotedblright{}, and Sanskrit \emph{dāru}, \textquotedblleft{}wood\textquotedblright{}; Trubachev adds Gothic \emph{triu}, so English \emph{tree} is a cousin, from PIE \emph{*dóru}, \textquotedblleft{}tree, wood\textquotedblright{}.
\end{cousinweb}

\begin{grammarlens}[title={What you've built}]
One more word for this chapter.
\end{grammarlens}

\subsection*{Your turn}
\begin{itemize}
\raggedright
  \item \emph{Say it:} \emph{dérevo}
  \item \emph{Say it:} \emph{dérevo}, once more
  \item \emph{Say it:} say \emph{dérevo}
  \item \emph{Recall:} say the Russian for a ticket, then the Russian for a map, then say \emph{dérevo} again
\end{itemize}

\begin{tcolorbox}[breakable,colback=teal!4,colframe=teal!35!black,title={Before you move on}]
What does \ru{дерево} mean? (\textquotedblleft{}A tree\textquotedblright{}.) Say it once more.
\end{tcolorbox}
\section[\texorpdfstring{\ru{цветок} (tsvetók)}{цветок (tsvetók)}]{\ru{цветок} (tsvetók) — a flower}

\label{lesson:RU-C55-flower}



\begin{quote}
Before the new one: say the Russian for a map, then the Russian for a tree.
\end{quote}

\subsection*{You'll want to know: \ru{цветок}}
\textbf{\ru{цветок}} (\emph{tsvetók}) — \textquotedblleft{}a flower\textquotedblright{}.

\begin{cousinweb}
\textbf{\ru{цветок}} is \emph{\ru{цвет}}, \textquotedblleft{}bloom; colour\textquotedblright{}, with the diminutive \emph{-\ru{ок}}. \emph{\ru{цвет}} $\leftarrow$ Proto-Slavic \emph{*květŭ}, \textquotedblleft{}blossom\textquotedblright{}, which Vasmer relates to Latvian \emph{kvitêt}, \textquotedblleft{}to shimmer\textquotedblright{}.
\end{cousinweb}

\begin{grammarlens}[title={What you've built}]
One more word for this chapter.
\end{grammarlens}

\subsection*{Your turn}
\begin{itemize}
\raggedright
  \item \emph{Say it:} \emph{tsvetók}
  \item \emph{Say it:} \emph{tsvetók}, once more
  \item \emph{Say it:} say \emph{tsvetók}
  \item \emph{Recall:} say the Russian for a map, then the Russian for a tree, then say \emph{tsvetók} again
\end{itemize}

\begin{tcolorbox}[breakable,colback=teal!4,colframe=teal!35!black,title={Before you move on}]
What does \ru{цветок} mean? (\textquotedblleft{}A flower\textquotedblright{}.) Say it once more.
\end{tcolorbox}
\section[\texorpdfstring{\ru{лист} (list)}{лист (list)}]{\ru{лист} (list) — a leaf; a sheet of paper}

\label{lesson:RU-C55-leaf}



\begin{quote}
Before the new one: say the Russian for a tree, then the Russian for a flower.
\end{quote}

\subsection*{You'll want to know: \ru{лист}}
\textbf{\ru{лист}} (\emph{list}) — \textquotedblleft{}a leaf; a sheet of paper\textquotedblright{}.

\begin{cousinweb}
\textbf{\ru{лист}} $\leftarrow$ Proto-Slavic \emph{*listŭ} — inherited, related to Lithuanian \emph{laĩškas}, \textquotedblleft{}leaf\textquotedblright{}, and Latvian \emph{laiska}, \textquotedblleft{}a leaf of flax\textquotedblright{} (Vasmer). Paper took the same word, as in English \emph{a leaf of a book}.
\end{cousinweb}

\begin{grammarlens}[title={What you've built}]
One more word for this chapter.
\end{grammarlens}

\subsection*{Your turn}
\begin{itemize}
\raggedright
  \item \emph{Say it:} \emph{list}
  \item \emph{Say it:} \emph{list}, once more
  \item \emph{Say it:} say \emph{list}
  \item \emph{Recall:} say the Russian for a tree, then the Russian for a flower, then say \emph{list} again
\end{itemize}

\begin{tcolorbox}[breakable,colback=teal!4,colframe=teal!35!black,title={Before you move on}]
What does \ru{лист} mean? (\textquotedblleft{}A leaf; a sheet of paper\textquotedblright{}.) Say it once more.
\end{tcolorbox}
\section[\texorpdfstring{\ru{трава} (travá)}{трава (travá)}]{\ru{трава} (travá) — grass}

\label{lesson:RU-C55-grass}



\begin{quote}
Before the new one: say the Russian for a flower, then the Russian for a leaf; a sheet of paper.
\end{quote}

\subsection*{You'll want to know: \ru{трава}}
\textbf{\ru{трава}} (\emph{travá}) — \textquotedblleft{}grass\textquotedblright{}.

\begin{cousinweb}
\textbf{\ru{трава}} $\leftarrow$ Proto-Slavic \emph{*trava} — inherited, from the root of Old Russian \emph{\ru{трути}}, \textquotedblleft{}to use up, to consume\textquotedblright{} (Vasmer): grass as the fodder that is eaten. The same root is in \emph{\ru{травить}}, \textquotedblleft{}to poison\textquotedblright{}.
\end{cousinweb}

\begin{grammarlens}[title={What you've built}]
One more word for this chapter.
\end{grammarlens}

\subsection*{Your turn}
\begin{itemize}
\raggedright
  \item \emph{Say it:} \emph{travá}
  \item \emph{Say it:} \emph{travá}, once more
  \item \emph{Say it:} say \emph{travá}
  \item \emph{Recall:} say the Russian for a flower, then the Russian for a leaf; a sheet of paper, then say \emph{travá} again
\end{itemize}

\begin{tcolorbox}[breakable,colback=teal!4,colframe=teal!35!black,title={Before you move on}]
What does \ru{трава} mean? (\textquotedblleft{}Grass\textquotedblright{}.) Say it once more.
\end{tcolorbox}
\section[\texorpdfstring{\ru{деревня} (derévnya)}{деревня (derévnya)}]{\ru{деревня} (derévnya) — a village}

\label{lesson:RU-C55-village}



\begin{quote}
Before the new one: say the Russian for a leaf; a sheet of paper, then the Russian for grass.
\end{quote}

\subsection*{You'll want to know: \ru{деревня}}
\textbf{\ru{деревня}} (\emph{derévnya}) — \textquotedblleft{}a village\textquotedblright{}.

\begin{cousinweb}
\textbf{\ru{деревня}} first meant \textquotedblleft{}a ploughed field\textquotedblright{}, then \textquotedblleft{}a farmstead with its land\textquotedblright{}, then \textquotedblleft{}a village\textquotedblright{} (Vasmer). It goes back to \emph{*dĭrvĭnja}, and Vasmer reports the usual link with Lithuanian \emph{dirvà}, \textquotedblleft{}a field\textquotedblright{}.
\end{cousinweb}

\begin{grammarlens}[title={What you've built}]
One more word for this chapter.
\end{grammarlens}

\subsection*{Your turn}
\begin{itemize}
\raggedright
  \item \emph{Say it:} \emph{derévnya}
  \item \emph{Say it:} \emph{derévnya}, once more
  \item \emph{Say it:} say \emph{derévnya}
  \item \emph{Recall:} say the Russian for a leaf; a sheet of paper, then the Russian for grass, then say \emph{derévnya} again
\end{itemize}

\begin{tcolorbox}[breakable,colback=teal!4,colframe=teal!35!black,title={Before you move on}]
What does \ru{деревня} mean? (\textquotedblleft{}A village\textquotedblright{}.) Say it once more.
\end{tcolorbox}
